% =====================================================================
% Conference draft (follow-up to "Model-Free PID Tuning by
% Step-Response Inspection"). Target: 6 pages.
%
% Class note: written for elsarticle.cls (present in paper/ of the
% project). If ADCHEM is chosen, port to ifacconf.cls -- the body
% transfers unchanged; only the front matter macros differ. Re-check
% the 6-page budget there (two-column, denser).
%
% Numbers policy: every number below is copied from the verified
% follow-up runs (code/asymptotics/*.py). Do not edit by hand;
% regenerate. Entries marked TODO-REGEN still need a regeneration
% script before submission.
% =====================================================================
\documentclass[preprint,12pt]{elsarticle}

\usepackage{amsmath,amssymb,amsthm}
\usepackage{graphicx}
\usepackage{booktabs}

\newtheorem{proposition}{Proposition}
\newtheorem{corollary}{Corollary}
\newtheorem{lemma}{Lemma}
\newtheorem{remark}{Remark}

\newcommand{\Nbar}{\bar{N}}
\newcommand{\Neps}{N_{\varepsilon}}

\begin{document}

\begin{frontmatter}

\title{Limit cycles and regularization in step-response
phase-portrait PID tuning}

% TODO: authors, affiliations, CRediT, funding -- same placeholders
% as in paper 1; fix both together before submission.
\author{Author One}
\author{Author Two}
\address{Placeholder affiliation}

\begin{abstract}
In a companion paper we introduced SPIN, a model-free PID tuning
method that inspects three phase portraits of the closed-loop step
response, counts a turn index in each, and adjusts the gains by a
fixed table of four moves. The method works, but the companion paper
does not say what the iteration converges \emph{to}. This paper
answers that question. First, we show that in log-gain coordinates
the four moves admit a unique vanishing convex combination, so the
only possible resting target is the \emph{triple point}: the gain
vector at which all three turn indices sit exactly on their
thresholds. When the triple point falls outside the gain box, the
target degenerates to a complementarity point, which explains a
boundary run reported in the companion paper. Second, we show that
the published turn index is a discontinuous function of the gains, so
the triple-point equation generally has no solution and
diminishing-step runs freeze. Third, we repair this with an
$\varepsilon$-regularized winding integral, after which
diminishing-step runs converge to the triple point. Finally, we
quantify the limit cycles of constant-step runs and compare step-size
schedules; a simple reversal-triggered halving schedule reaches
$10^{-4}$ accuracy in under a hundred iterations.
\end{abstract}

\begin{keyword}
PID control \sep model-free tuning \sep winding number \sep
switched systems \sep limit cycles
\end{keyword}

\end{frontmatter}

% =====================================================================
\section{Introduction}
\label{sec:intro}

Most PID tuning methods fit a model first and compute gains second
\cite{ZN1942,AstromHagglund2004}. In a companion paper
\cite{Paper1} we took a different route: tune directly from what the
step response \emph{looks like}. The method, SPIN (Step-response
Phase-portrait INspection), draws three phase portraits from one
closed-loop step response, counts how many times each curve turns
around the origin --- the \emph{turn index} $N_k$ of band $k$ --- and
compares each count to a fixed threshold. A table of four moves, the
\emph{triangular rule}, then scales the gains up or down. No plant
model is identified at any point.

The companion paper demonstrates that the rule detunes aggressive
gains, recovers from failed loops, and lands in a feasible region on
a battery of standard plants. What it does not do is characterize
the iteration's asymptotic behavior. Three questions were left open.
What point, if any, does the rule aim at? Why do runs with a
shrinking step size freeze at points that depend on the schedule
rather than on the plant? And which step-size schedules actually
converge?

This paper answers all three. The answers turn out to be connected.
The rule aims at a single well-defined target, the \emph{triple
point}, at which all three turn indices sit exactly on their
thresholds (Section~\ref{sec:triple}). Diminishing-step runs freeze
because the published turn index is discontinuous in the gains, so
the triple-point equation generally has no solution
(Section~\ref{sec:obstruction}). Replacing the hard count with an
$\varepsilon$-regularized winding integral removes the discontinuity,
and with it the freezing (Section~\ref{sec:regularized}). With a
continuous index in hand, the limit cycles of constant-step runs and
the convergence of diminishing-step schedules can be measured
cleanly (Section~\ref{sec:schedules}). Section~\ref{sec:structure}
examines the local structure at the target and identifies the route
a convergence proof must take.

It is worth stating what kind of object we are analyzing. Iterative
model-free tuning is not new: iterative feedback tuning estimates
the gradient of a performance criterion from closed-loop experiments
and converges, with an appropriately shrinking step size, to a local
minimum of that criterion \cite{Hjalmarsson1998}; extremum seeking
does the same with sinusoidal probing \cite{Killingsworth2006}. The
triangular rule is different in kind. It is a \emph{sign-based}
update: each iteration compares three integers (later: three real
numbers) to three thresholds and picks one of four fixed moves.
There is no cost function and no gradient, so the limit object is
not a stationary point of anything. It is a point where a switching
law balances --- and the right mathematical frame is that of switched
systems and their averaged dynamics, not descent.

% =====================================================================
\section{The rule as a switched system in log-gain coordinates}
\label{sec:rule}

We recall only what is needed; the companion paper \cite{Paper1}
defines the portraits, the index, and the rule in full.

One closed-loop step experiment yields three curves (band $0$, $1$,
$2$), and for each curve a turn index $N_k \ge 0$ measuring how far
the curve winds around the origin of its portrait. Small indices are
good. Each band has a fixed threshold, and in this paper the
thresholds are
\begin{equation}
\Nbar = (\Nbar_0, \Nbar_1, \Nbar_2) = (0.5,\; 0.75,\; 1.0),
\label{eq:thresholds}
\end{equation}
the same values used throughout the companion paper. The triangular
rule reads the three comparisons $N_k \gtrless \Nbar_k$ and applies
one of four moves to the gain vector $K = (K_p, K_i, K_d)$: if all
three bands are below threshold, all gains are scaled up; otherwise
the lowest offending band is cut by scaling gains down in a
band-specific pattern. Every move multiplies each gain by
$(1+\beta)^{\pm 1}$ or leaves it alone, where $\beta > 0$ is the
step size. Gains are clipped to a box, here $[0.01, 10]$ for each
gain.

Because the moves are multiplicative, the natural coordinates are
logarithmic. Let $\theta = (\log K_p, \log K_i, \log K_d)$ and
$h = \log(1+\beta)$. The four moves become translations
\begin{equation}
\begin{aligned}
e   &= (+1, +1, +1) &\quad&\text{(expand: all bands below threshold),}\\
d_0 &= (-1, \phantom{+}0, \phantom{+}0) &&\text{(cut band 0),}\\
d_1 &= (+1, -1, \phantom{+}0) &&\text{(cut band 1),}\\
d_2 &= (+1, +1, -1) &&\text{(cut band 2),}
\end{aligned}
\label{eq:moves}
\end{equation}
each scaled by $h$. The iteration is therefore a switched dynamical
system on $\theta$: a partition of gain space into four regions
(determined by the signs of $N_k(\theta) - \Nbar_k$), with a
constant velocity assigned to each region. Figure~\ref{fig:moves}
shows the geometry.

% TODO-FIG: figures/moves_geometry.png -- the four move vectors in
% log-gain space; generate with code/asymptotics/plot_moves.py
\begin{figure}[t]
\centering
\fbox{\parbox[c][4.5cm][c]{0.8\linewidth}{\centering
TODO: move geometry figure}}
\caption{The four moves \eqref{eq:moves} of the triangular rule in
log-gain coordinates. Each iteration applies exactly one of the four
translations, selected by comparing the three turn indices to their
thresholds.}
\label{fig:moves}
\end{figure}

Two properties of this system drive everything that follows. First,
the velocities are constant within each region, so nothing can
converge unless the state reaches a place where the regions meet and
the moves can balance on average. Second, the switching surfaces are
the level sets $N_k(\theta) = \Nbar_k$, so their regularity --- or
lack of it --- decides whether such a balance point exists at all.

% =====================================================================
\section{The target: a triple point}
\label{sec:triple}

Suppose a trajectory of the rule stays bounded and does not drift.
Over a long window, each move $v \in \{e, d_0, d_1, d_2\}$ is applied
some fraction $w_v \ge 0$ of the time, with $\sum_v w_v = 1$, and the
average displacement per iteration is $h \sum_v w_v v$. A
non-drifting trajectory requires this average to vanish.

\begin{proposition}[Unique balance]
\label{prop:balance}
The only convex combination of the moves \eqref{eq:moves} that
vanishes is
\begin{equation}
\tfrac{1}{8}\, e + \tfrac{1}{2}\, d_0 + \tfrac{1}{4}\, d_1
+ \tfrac{1}{8}\, d_2 = 0 .
\label{eq:weights}
\end{equation}
\end{proposition}

\begin{proof}
The third coordinate of \eqref{eq:moves} is nonzero only for $e$
($+1$) and $d_2$ ($-1$), so $w_e = w_{d_2}$. The second coordinate
then forces $w_{d_1} = w_e + w_{d_2} = 2 w_e$, and the first forces
$w_{d_0} = w_e + w_{d_1} + w_{d_2} = 4 w_e$. Normalizing gives
\eqref{eq:weights}.
\end{proof}

All four weights in \eqref{eq:weights} are strictly positive. This
has a sharp consequence: a bounded, non-drifting trajectory must
keep visiting all four regions, which is only possible if it hovers
where all three switching surfaces meet.

\begin{corollary}[Triple point]
\label{cor:triple}
The only possible resting target of the rule in the interior of the
gain box is a point $\theta^*$ with
\begin{equation}
N_k(\theta^*) = \Nbar_k, \qquad k = 0, 1, 2,
\label{eq:triplepoint}
\end{equation}
i.e.\ all three turn indices exactly on threshold. Near
$\theta^*$ the moves are applied with limiting frequencies
$(w_e, w_{d_0}, w_{d_1}, w_{d_2})
= (\tfrac18, \tfrac12, \tfrac14, \tfrac18)$.
\end{corollary}

The corollary is testable, and it holds with striking precision.
On plants P1, P3 and P4 of the companion battery, long constant-step
runs settle into a periodic orbit around a single point, and the
measured duty cycle of the four moves matches
$(\tfrac18, \tfrac12, \tfrac14, \tfrac18)$ to three decimal places,
at both $\beta = 0.10$ and $\beta = 0.03$. The terminal orbit itself
is a period-8 cycle in which the moves appear in a fixed
``ruler'' pattern --- the shortest sequence realizing the weights
$(1, 4, 2, 1)/8$. Figure~\ref{fig:ruler} shows one period.
Table~\ref{tab:duty} summarizes the duty-cycle measurements.

% TODO-FIG: figures/ruler_orbit.png -- one period of the terminal
% orbit around the triple point on P1, moves color-coded.
\begin{figure}[t]
\centering
\fbox{\parbox[c][4.5cm][c]{0.8\linewidth}{\centering
TODO: period-8 terminal orbit figure}}
\caption{Terminal orbit of a constant-step run on P1
($\beta = 0.03$): a period-8 cycle around the triple point
$\theta^*$. The move sequence realizes the balance weights
\eqref{eq:weights} exactly once per period.}
\label{fig:ruler}
\end{figure}

\begin{table}[t]
\centering
\caption{Measured duty cycle of the four moves over the terminal
orbit, against the prediction of Corollary~\ref{cor:triple}.
% TODO-REGEN: measured columns from code/asymptotics -- add a
% reproduce script before submission; values matched prediction to
% three decimals on P1, P3, P4 at both step sizes.
}
\label{tab:duty}
\begin{tabular}{lcccc}
\toprule
 & $w_e$ & $w_{d_0}$ & $w_{d_1}$ & $w_{d_2}$ \\
\midrule
Predicted & $0.125$ & $0.500$ & $0.250$ & $0.125$ \\
P1, $\beta{=}0.10$ & \multicolumn{4}{c}{TODO-REGEN} \\
P1, $\beta{=}0.03$ & \multicolumn{4}{c}{TODO-REGEN} \\
P3, P4 (both $\beta$) & \multicolumn{4}{c}{TODO-REGEN} \\
\bottomrule
\end{tabular}
\end{table}

\paragraph{Boundary case} If $\theta^*$ falls outside the gain box,
condition \eqref{eq:triplepoint} cannot be met and the target
degenerates: the trajectory settles at a point where, for each band,
either the index sits on its threshold \emph{or} the associated gain
sits on a bound --- a complementarity point. The companion paper
contains an instance without naming it as such: on plant P2, the
demonstration run ends with $K_d$ resting on the $0.01$ floor and
terminal indices $N = (0.11,\, 1.04,\, 1.14)$ at gains
$(0.392,\, 0.025,\, 0.01)$. Two indices hover at their thresholds
while the third constraint is replaced by the active gain bound.
What the companion paper describes as ``the rule removing derivative
action on a delay-dominant plant'' is exactly the complementarity
target predicted here.

% =====================================================================
\section{The obstruction: a discontinuous count}
\label{sec:obstruction}

Corollary~\ref{cor:triple} presumes that equation
\eqref{eq:triplepoint} has a solution. With the turn index as
published, it generally does not.

The published index is a hard count: the portrait curve is followed,
its winding about the origin is accumulated, and the accumulation is
truncated when the curve enters a small disc of radius
$\varepsilon$ about the origin for the last time (a $\delta$-guard
suppresses spurious re-entries). The truncation is the problem. The
time of \emph{last entry} into the disc is not a continuous
functional of the curve: an arbitrarily small deformation can make
the curve graze the disc boundary and either gain or lose a whole
excursion.

The effect is easy to produce. On the battery plants, a gain change
of $0.05\,\%$ --- far below any step the rule ever takes --- moves one
turn index by exactly $1.000$. Lengthening the simulation window
does not help: the jump is caused by the truncation rule, not by the
window. Figure~\ref{fig:jump} shows the index as a function of a
single gain through such a jump.

% TODO-FIG: figures/index_jump.png -- N_k vs. one gain across the
% discontinuity; hard count vs. regularized index overlaid.
\begin{figure}[t]
\centering
\fbox{\parbox[c][4.5cm][c]{0.8\linewidth}{\centering
TODO: discontinuity figure (hard count vs.\ $\Neps$)}}
\caption{One turn index as a single gain is varied by $0.05\,\%$.
The hard-truncated count (dashed) jumps by exactly one; the
regularized index $\Neps$ of Section~\ref{sec:regularized} (solid)
moves by $7 \times 10^{-4}$.}
\label{fig:jump}
\end{figure}

The consequence for the asymptotics is immediate. The switching
surfaces $N_k(\theta) = \Nbar_k$ are, in general, not surfaces at
all: the index jumps across them in integer steps, and the exact
threshold value is attained on a measure-zero set or not at all. The
triple-point equation \eqref{eq:triplepoint} then has no solution,
and a run whose step size shrinks to zero must freeze at whatever
point its remaining travel budget allows --- a point that depends on
the schedule, not the plant. This is precisely the freezing observed
in practice, and it is why the companion paper used a constant step
size throughout.

For a constant step the discontinuity is harmless: the rule only
reads the \emph{sign} of $N_k - \Nbar_k$, and a jump of one merely
relocates the switching surface slightly. This is why the companion
paper's results stand as published. The discontinuity only becomes
fatal when one asks the iteration to converge to a point, which is
what this paper asks.

% =====================================================================
\section{The fix: a regularized turn index}
\label{sec:regularized}

The repair is to stop truncating and start weighting. Recall that
for a closed curve $(x(t), y(t))$ avoiding the origin, the winding
number is the integral $\frac{1}{2\pi} \oint
\frac{x\,dy - y\,dx}{x^2 + y^2}$. The hard count approximates this
integral while excising the part of the curve inside the
$\varepsilon$-disc. Instead of excising, we damp:

\begin{equation}
\Neps \;=\; \frac{1}{2\pi} \oint
\frac{x\,dy - y\,dx}{\,x^2 + y^2 + \varepsilon^2\,}.
\label{eq:neps}
\end{equation}

Writing $\rho^2 = x^2 + y^2$, the integrand of \eqref{eq:neps} is
the ordinary winding integrand multiplied by the weight
$\rho^2 / (\rho^2 + \varepsilon^2)$: turning far from the origin
counts fully, turning at radius $\varepsilon$ counts half, and
turning deep inside the disc counts almost nothing. The disc's sharp
boundary --- the source of the discontinuity --- is gone, and with it
the need to decide when the curve has entered the disc ``for the
last time''.

Three properties matter in practice.

\emph{Continuity.} The integrand in \eqref{eq:neps} is smooth in
$(x, y)$ everywhere, including the origin, so $\Neps$ is a
continuous functional of the curve, hence a continuous function of
the gains wherever the step response depends continuously on them.
On the perturbation of Figure~\ref{fig:jump}, the change in $\Neps$
is $7 \times 10^{-4}$ where the hard count jumps by $1.000$.

\emph{Window independence.} The tail of a settled response
contributes radius $\rho \to 0$, where the weight
$\rho^2/(\rho^2+\varepsilon^2)$ vanishes quadratically. The value of
$\Neps$ therefore becomes exactly independent of the simulation
window once the response has settled, and the $\delta$-guard of the
hard count is unnecessary. Nothing about the experiment or the
portraits changes; only the counting does.

\emph{Interchangeability.} Away from the switching surfaces,
$\Neps$ and the hard count agree to within a fraction that vanishes
as $\varepsilon \to 0$, so every sign comparison the published rule
makes is unchanged. All results of the companion paper hold verbatim
with $\Neps$ in place of the hard count. We keep the thresholds
\eqref{eq:thresholds} and the value of $\varepsilon$ exactly as
published.
% TODO: one sentence stating the published epsilon value here, for
% self-containedness (copy from paper 1, Sec. on the index).

With $\Neps$, the switching surfaces are genuine level sets, the
triple-point equation \eqref{eq:triplepoint} has a solution, and
diminishing-step runs stop freezing. On the three interior-target
plants, runs with vanishing steps now converge to the triple point
with final log-gain errors
\begin{equation*}
\text{P1: } 0.010, \qquad \text{P3: } 0.017, \qquad \text{P4: } 0.041,
\end{equation*}
where the error is $\| \theta_{\mathrm{final}} - \theta^* \|$ and
$\theta^*$ is located independently by direct search on
\eqref{eq:triplepoint}. The freezing of Section~\ref{sec:obstruction}
is gone; what remains is the ordinary trade-off between step
schedule and speed, which we now quantify.

% =====================================================================
\section{Step-size schedules and limit cycles}
\label{sec:schedules}

All runs in this section start far from the target --- gains detuned
by a factor of $16$ --- and use the regularized index. We compare a
constant step against diminishing schedules $\beta_n \propto
n^{-1/2}$, $n^{-3/4}$, $n^{-1}$, and against a simple adaptive
schedule. Table~\ref{tab:schedules} and Figure~\ref{fig:schedules}
summarize the outcome; errors are final log-gain distances to
$\theta^*$.

\begin{table}[t]
\centering
\caption{Final log-gain error $\|\theta - \theta^*\|$ for different
step-size schedules, far start ($16\times$ detuned), regularized
index.
% TODO-REGEN: confirm whether these are per-plant or averaged over
% P1/P3/P4 and add the reproduce script; values from
% code/asymptotics runs of July 2026.
}
\label{tab:schedules}
\begin{tabular}{lll}
\toprule
Schedule & Final error & Behaviour \\
\midrule
$\beta$ constant           & $0.074$ & limit cycle, $O(\beta)$ radius \\
$\beta_n \propto n^{-1/2}$ & $0.010$ & converges \\
$\beta_n \propto n^{-3/4}$ & $0.001$ & converges \\
$\beta_n \propto n^{-1}$   & $0.197$ & freezes short (Lemma~\ref{lem:travel}) \\
reversal-halving           & $10^{-4}$ & converges, $36$--$98$ iterations \\
\bottomrule
\end{tabular}
\end{table}

% TODO-FIG: figures/schedule_convergence.png -- error vs. iteration
% for the five schedules, log scale.
\begin{figure}[t]
\centering
\fbox{\parbox[c][4.5cm][c]{0.8\linewidth}{\centering
TODO: convergence curves figure}}
\caption{Distance to the triple point versus iteration for the
schedules of Table~\ref{tab:schedules} (log scale). The constant
step settles into the period-8 orbit of Figure~\ref{fig:ruler}; the
$n^{-1}$ schedule stalls before reaching the target.}
\label{fig:schedules}
\end{figure}

\paragraph{Constant step} A constant $\beta$ cannot converge to a
point: near $\theta^*$ the state is carried around the period-8
orbit of Section~\ref{sec:triple}, whose radius scales with the step
$h = \log(1+\beta)$. The measured final error of $0.074$ is this
orbit radius, not a failure of the rule. For plant-floor use, where
gains within a few percent of ideal are indistinguishable, the
constant step remains a perfectly reasonable choice.

\paragraph{Diminishing steps} The schedules $n^{-1/2}$ and
$n^{-3/4}$ converge, and faster decay gives a smaller final error.
The classical schedule $\beta_n \propto n^{-1}$, however, freezes
well short of the target, at error $0.197$. The reason is a travel
budget, not noise:

\begin{lemma}[Logarithmic travel]
\label{lem:travel}
Under $\beta_n = c/n$, the total distance the iteration can cover in
$n$ steps is $\sum_{m \le n} \log(1 + c/m) = O(\log n)$. A run
started at distance $D$ from the target cannot reach it before
$n \approx e^{D/c}$ iterations even if every step points the right
way.
\end{lemma}

With a far start ($D \approx \log 16$ per gain), the required $n$
is astronomically large, and the run stalls at whatever distance the
first few hundred iterations reach. This is the same summability
boundary known from stochastic approximation \cite{RobbinsMonro1951},
met here from the travel side rather than the noise side: our
iteration is deterministic, and $n^{-1}$ fails not because averaging
is too slow but because the path itself is too short.

\paragraph{Reversal-triggered halving} The best schedule we found
needs no exponent at all: hold $\beta$ constant until every band has
reversed its move direction at least once --- the sign pattern
$(N_k - \Nbar_k)$ has flipped in each coordinate --- then halve
$\beta$ and repeat. Reversals in all bands indicate that the
trajectory has begun to orbit rather than travel, so the step is no
longer limiting progress and can be safely cut. This schedule
reaches error $10^{-4}$ in $36$--$98$ iterations across the
interior-target plants. Readers from the learning literature will
recognize the mechanism: it is a global cousin of the sign-based
Rprop update with reversal-triggered step adaptation
\cite{Riedmiller1993}, arising here naturally because the triangular
rule is itself sign-based.

% =====================================================================
\section{Local structure and the route to a proof}
\label{sec:structure}

Is the triple point locally unique, and is convergence to it
provable? The first question we can answer numerically; the second
we can at least point in the right direction.

Central-difference Jacobians $\partial N_j / \partial \theta_i$ of
the regularized index at $\theta^*$ are nonsingular on all three
interior-target plants (determinants $1.2$, $5.9$, $4.3$ on P1, P3,
P4), so $\theta^*$ is locally unique in each case.
% TODO-REGEN: units/normalization of the determinants; add to the
% reproduce script.

The Jacobians are, however, decidedly \emph{not} diagonal, and not
even diagonally dominant. On P1, the proportional gain moves the
band-2 index harder than the derivative gain does; on P4, the
own-gain entry of band~1 is slightly negative. The intuitive picture
behind the triangular rule --- each move mainly fixes its own band ---
is therefore false at the target: a per-move descent argument, in
which each cut is shown to reduce its own band's index, fails on two
of the three plants. Yet all runs converge.

The resolution is that convergence is a property of the
\emph{averaged} dynamics, not of any single move. Near $\theta^*$
the iteration chatters between the four regions, and its mean motion
follows the convex combination of the region velocities weighted by
local residence times --- the Filippov regularization of the switched
system \cite{Filippov1988}. It is this averaged vector field, not
the individual rows of the Jacobian, that points toward $\theta^*$;
the duty-cycle measurements of Section~\ref{sec:triple} are direct
evidence, since the observed residence times are exactly the balance
weights \eqref{eq:weights}. A convergence proof should therefore
establish stability of $\theta^*$ for the Filippov dynamics and
conclude for the discrete iteration by a standard
averaging/perturbation step. We record this as the open problem of
the paper. A practical corollary is worth stating: acceleration
schemes that treat the bands independently (per-band step sizes,
per-band momentum) implicitly assume the diagonal structure that is
absent here and should be expected to misbehave; momentum applied to
the full move vector, which smooths the chattering the way the
Filippov average does, is the variant worth trying.

% =====================================================================
\section{Conclusions}
\label{sec:conclusions}

The triangular rule of the companion paper is a switched dynamical
system whose only interior resting target is the triple point: all
three turn indices exactly on threshold, approached through a
period-8 orbit whose move frequencies are fixed by a unique convex
balance. When the triple point leaves the gain box, the target
degenerates to a complementarity point, which the companion paper's
P2 run realizes. The published hard-truncated count is discontinuous
in the gains and therefore cannot support point convergence; the
$\varepsilon$-regularized winding integral \eqref{eq:neps} restores
continuity at no cost to any published result. With it,
diminishing-step runs converge, the $n^{-1}$ schedule is shown to
fail for a travel-budget reason, and a reversal-triggered halving
schedule reaches $10^{-4}$ accuracy in under a hundred iterations.
The convergence proof itself runs through the averaged Filippov
dynamics and is left open, together with the momentum variant it
suggests.

% TODO: Acknowledgments / funding -- align with paper 1 placeholders.

% =====================================================================
\begin{thebibliography}{10}

\bibitem{Paper1}
Authors, Model-free PID tuning by step-response inspection,
submitted to J. Process Control, 2026.
% TODO: update once paper 1 has a status/DOI.

\bibitem{ZN1942}
J.~G. Ziegler, N.~B. Nichols, Optimum settings for automatic
controllers, Trans. ASME 64 (1942) 759--768.

\bibitem{AstromHagglund2004}
K.~J. {\AA}str{\"o}m, T. H{\"a}gglund, Revisiting the
Ziegler--Nichols step response method for PID control,
J. Process Control 14 (2004) 635--650.

\bibitem{Hjalmarsson1998}
H. Hjalmarsson, M. Gevers, S. Gunnarsson, O. Lequin, Iterative
feedback tuning: theory and applications, IEEE Control Syst. Mag.
18 (4) (1998) 26--41.

\bibitem{Killingsworth2006}
N.~J. Killingsworth, M. Krsti{\'c}, PID tuning using extremum
seeking, IEEE Control Syst. Mag. 26 (1) (2006) 70--79.

\bibitem{RobbinsMonro1951}
H. Robbins, S. Monro, A stochastic approximation method,
Ann. Math. Statist. 22 (1951) 400--407.

\bibitem{Riedmiller1993}
M. Riedmiller, H. Braun, A direct adaptive method for faster
backpropagation learning: the RPROP algorithm, in: Proc. IEEE Int.
Conf. Neural Networks, 1993, pp. 586--591.

\bibitem{Filippov1988}
A.~F. Filippov, Differential Equations with Discontinuous
Righthand Sides, Kluwer, 1988.

\end{thebibliography}

\end{document}
